% ============================================================
% f2014_beijing_arts_q18_histogram.tex
% 2014 北京卷（文科）第 18 题：100 名学生一周课外阅读时间的频率分布直方图
% 来源 PDF：普通高考/2014/2014北京文.pdf
% 原页：第 2 页（A3），印面 864，第 18 题频数分布表右侧
% 组距 2，共 9 组；柱高 = 频数/200，其中第 3 组顶端标 a、第 5 组顶端标 b。
% 横轴每 2 小时 1.5cm，纵轴每 0.025 为 0.75cm（b=0.125 对应 3.75cm）。
% ============================================================
\begin{tikzpicture}[>=Stealth, line width=0.6pt, font=\small,
  x=0.75cm, y=0.75cm]
  % 坐标轴
  \draw[->] (0,0) -- (0,5.75);
  \draw[->] (0,0) -- (20.6,0);

  % 纵轴标题（频率/组距）
  \draw (0.05,6.40) -- (1.55,6.40);
  \node[above] at (0.80,6.40) {频率};
  \node[below] at (0.80,6.40) {组距};

  % 横轴刻度
  \node[below left] at (0,0) {$0$};
  \foreach \x in {2,4,6,8,10,12,14,16,18} {
    \draw (\x,0) -- (\x,-0.14);
    \node[below] at (\x,-0.02) {$\x$};
  }

  % 九个矩形：频数 6,8,17,22,25,12,6,2,2
  \foreach \i/\n in {0/6,1/8,2/17,3/22,4/25,5/12,6/6,7/2,8/2} {
    \draw (\i*2,0) rectangle (\i*2+2,\n*0.16);
  }

  % a、b 的参考虚线与标注
  \draw[dashed, line width=0.45pt] (0,2.72) -- (6,2.72);
  \draw[dashed, line width=0.45pt] (0,4.00) -- (10,4.00);
  \node[left] at (0,2.72) {$a$};
  \node[left] at (0,4.00) {$b$};

  % 横轴名称
  \node[above] at (19.1,0.12) {阅读时间};
\end{tikzpicture}
